%% ma: L12-B14-0146
%% lop: 12
%% chuong: 5
%% bai: 14
%% dang: 12.14.29
%% loai: TLN
%% muc: TH
%% dap_an: "-2"
%% nguon: "Ảnh giáo viên gửi 25/11/2026 (tài liệu ôn tập Phương trình mặt phẳng), Phần III Câu 3 (điểm M(2;1;-3), mặt phẳng chứa trục Ox)"
%% kien_thuc: "Bài 14 (mặt phẳng chứa trục tọa độ)"
%% trang_thai: chinh-thuc
\begin{ex}
Trong không gian $Oxyz$, cho điểm $M(2;1;-3)$. Phương trình mặt phẳng $(\alpha)$ đi qua $M$ và chứa trục $Ox$ có dạng $ax+3y+cz+d=0$. Tính $a-2c+d$.
\shortans{-2}
\loigiai{Mặt phẳng chứa $Ox$ đi qua $O$ và nhận $\vec i=(1;0;0)$ làm vectơ chỉ phương nên $a=0$ và $d=0$: $3y+cz=0$. Vì $M\in(\alpha)$: $3\cdot1+c\cdot(-3)=0\Rightarrow c=1$. Vậy $a-2c+d=0-2+0=-2$.}
\end{ex}
